\documentclass{article}

\usepackage[english]{babel}

\usepackage[letterpaper,top=2cm,bottom=2cm,left=3cm,right=3cm,marginparwidth=1.75cm]{geometry}

% Useful packages
\usepackage{amsmath}
\usepackage{graphicx}
\usepackage[colorlinks=true, allcolors=blue]{hyperref}

\title{Eigenvalues and Eigenvectors}
\author{Peter Zhao}

\begin{document}
\date{}
\maketitle

\section*{Acknowledgments}
I will be using various references pulled from different textbooks for a complete guide/notes to Eigenvalues. Here are some main textbooks of interest: 

\begin{itemize}
\item Gilbert Strang, \textit{Linear Algebra and Its Applications}, 4th Edition, MIT Press.

\item David C. Lay, Steven R. Lay, and Judi J. McDonald,
\textit{Linear Algebra and Its Applications}, 5th Edition,
Pearson, 2016.

\item Stephen H. Friedberg, Arnold J. Insel, and Lawrence E. Spence.
\textit{Linear Algebra}. 4th Edition. Pearson, 2003.


\end{itemize}

\section{Introduction}
\subsection{Background}
The topic of eigenvalues and eigenvectors is often taught in a way that feels abstract. 
I feel like students are usually given formulas and procedures without first understanding why these
objects are useful.

Eigenvalues arise when studying linear transformations. A linear transformation is a rule/function
that changes vectors in a space while preserving straight lines and the origin. Geometrically,
this transformation can stretch, compress, rotate, or shear the space.

To describe positions in a vector space, we use a set of vectors called a basis. A basis is
a collection of vectors that allows us to express every vector in the space as a linear
combination of those basis vectors.

When a matrix transformation acts on most vectors, it both changes their length and rotates
their direction. However, there are sometimes special vectors whose direction does not
change under the transformation. These vectors may only be stretched or compressed.

Such vectors are called eigenvectors.

The amount by which the transformation stretches or compresses an eigenvector is called
the eigenvalue.

Therefore, the eigenvalue problem asks us to find the vectors whose direction is preserved
by a linear transformation and determine how much those vectors are scaled.
Geometrically, eigenvectors are directions in space that remain unchanged by the transformation; only their length is scaled.

\subsection{Eigenvalue Problem Across Disciplines}

The eigenvalue problem appears in many areas of science and mathematics
and may be written using different notations:

\[
Av = \lambda v
\]

\[
T(v) = \lambda v
\]

\[
A|v\rangle = \lambda |v\rangle
\]

\[
(A - \lambda I)v = 0
\]

\[
\mathcal{L} f = \lambda f
\]

They all represent the same eigenvalue problem. 
Usually, we first encounter this concept in linear algebra written in matrix form as

\[
A\mathbf{v} = \lambda \mathbf{v}.
\]

Here $A$ is a square matrix, $\mathbf{v}$ is a nonzero vector called an eigenvector,
and $\lambda$ is the corresponding eigenvalue.

But interestingly, the eigenvalue problem does not only have to be in matrix form.
More generally, it appears when anything "linear".

For example, consider the differential linear operator

\[
\frac{d}{dt}.
\]

If we apply this operator to the function $e^{t}$, we obtain

\[
\frac{d}{dt} e^{t} = e^{t}.
\]

This can be interpreted as an eigenvalue problem:

\[
\frac{d}{dt} e^{t} = 1 \cdot e^{t}.
\]

Here $e^{t}$ plays the role of an eigenfunction, and the eigenvalue is $\lambda = 1$.

\subsection{A Simple Graph Ranking Example}

Before studying the mathematical theory in detail or solving it I feel it is helpful to
see an example of how eigenvalues appear in a real algorithm.


For example, a graph can be used to represent many systems in computer science,
such as webpages and hyperlinks.

Suppose each node in a graph represents a webpage, and a directed edge
represents a link from one page to another. The idea behind ranking
algorithms such as PageRank is that a page is important if important
pages link to it.

We begin by assigning every page the same score. Then we repeatedly
update the scores according to the links in the graph.

For example, consider three pages:

\begin{itemize}
\item Page $0$ links to Page $1$
\item Page $1$ links to Page $2$
\item Page $2$ links to Page $0$ and Page $1$
\end{itemize}

This structure can be represented by the adjacency matrix

\[
A =
\begin{bmatrix}
0 & 1 & 0 \\
0 & 0 & 1 \\
1 & 1 & 0
\end{bmatrix}
\]

Each row represents outgoing links from a page.

We start with an initial score vector

\[
\mathbf{s}_0 =
\begin{bmatrix}
1 \\
1 \\
1
\end{bmatrix}.
\]

At each step, the score of a page is distributed equally to the
pages it links to. After repeating this update many times, the
scores begin to stabilize.

The resulting stable score vector satisfies an equation of the form

\[
A\mathbf{s} = \lambda \mathbf{s}.
\]

This means that the final ranking vector $\mathbf{s}$ behaves like
an eigenvector of the matrix describing the graph.

Thus eigenvalues and eigenvectors help describe the long-term behavior
of ranking algorithms used in computer science.

If this example feels confusing, that is okay! The goal is to show that eigenvalues already appear in practical algorithms used in useful applications.

\section{Solving the Eigenvalue Problem}

Now that we understand what eigenvalues and eigenvectors represent,
we would like to compute them.

Recall the eigenvalue equation

\[
A\mathbf{v} = \lambda \mathbf{v}.
\]

We can rewrite this equation as

\[
A\mathbf{v} - \lambda \mathbf{v} = 0.
\]

Factoring out the vector $\mathbf{v}$ gives

\[
(A - \lambda I)\mathbf{v} = 0.
\]

Notice that $\lambda \mathbf{v}$ is equivalent to multiplying the vector
$\mathbf{v}$ by the scalar $\lambda$. In order to subtract this term from
$A\mathbf{v}$ using matrix notation, we introduce the identity matrix $I$. 
must be singular. Therefore we require

\[
\det(A - \lambda I) = 0.
\]

This equation is called the \textbf{characteristic equation}.
Solving this equation gives the eigenvalues of the matrix.





\end{document}
